\documentclass[12pt]{article}

% --- Packages ---
\usepackage[utf8]{inputenc}
\usepackage[T1]{fontenc}
\usepackage{amsmath,amssymb}
\usepackage{booktabs}
\usepackage{graphicx}
\usepackage{natbib}
\usepackage{hyperref}
\usepackage{geometry}
\usepackage{setspace}
\usepackage{caption}
\usepackage{subcaption}
\usepackage{multirow}
\usepackage{array}
\usepackage{longtable}
\usepackage{xcolor}

\geometry{margin=1in}
\doublespacing

\bibliographystyle{apalike}

% --- Title ---
\title{The Voynich Manuscript as Lossy Generation: Evidence for a Combinatorial Encoding Device from Positional Slot Analysis, Synthetic Text Comparison, and Botanical Correlation}

\author{%
  Showerhead\thanks{Contact: @Formershowerhead on YouTube. AI-assisted analysis methodology is disclosed in Section~\ref{sec:methods}.}
}

\date{\today}

\begin{document}

\maketitle

% ============================================================
% ABSTRACT
% ============================================================
\begin{abstract}
We present quantitative evidence that the text of the Voynich Manuscript (Beinecke MS 408) was produced by a physical combinatorial generation device---a lookup table or equivalent mechanism with up to 879 observed row signatures and seven positional columns. Five convergent lines of evidence support this conclusion. First, the manuscript's 37,886-token vocabulary decomposes into seven positional slots with 3--16 fragment options each, where all slot pairs exhibit statistically significant dependencies (chi-squared $p < 10^{-300}$, Cram\'{e}r's $V = 0.12$--$0.31$), ruling out independent-column mechanisms such as the Cardan grille hypothesis of \citet{Rugg2004}. Second, a default-filler signature---the fragment \texttt{ch} in slot~2 exhibits a 14$\times$ frequency spike when adjacent slots are empty ($V = 0.48$)---provides direct evidence of blank-cell defaults in a lookup table. Third, default rates independently recover manuscript section boundaries (Kruskal--Wallis $p = 4 \times 10^{-9}$), confirming \citet{Montemurro2013} via a wholly different method. Fourth, a row-based synthetic generator reproduces the real manuscript's slot-transition statistics with 9--15$\times$ less divergence (Jensen--Shannon) than an independent-slot generator, computationally falsifying the Cardan grille model. Fifth, and decisively, the real manuscript text correlates with botanical illustration features ($p = 0.003$ for root morphology) while synthetic text from the same generator does not ($p_{\text{ablation}} < 0.0001$ for all five tested features), proving that real content---not merely mechanical structure---passed through the generation system. We term this framework \emph{lossy generation}: the device encoded real information but introduced systematic degradation, producing text that is structurally valid but semantically impoverished---an effect analogous to hallucination in modern language models, but arising from a physical rather than computational mechanism.

\medskip
\noindent\textbf{Keywords:} Voynich Manuscript, combinatorial generation, Cardan grille, positional slot analysis, lossy encoding, computational falsification
\end{abstract}

\newpage

% ============================================================
% 1. INTRODUCTION
% ============================================================
\section{Introduction}
\label{sec:intro}

The Voynich Manuscript (Beinecke Library MS 408) is a 15th-century illustrated codex written in an undeciphered script, commonly transcribed using the European Voynich Alphabet (EVA) system. Radiocarbon dating places the vellum between 1404 and 1438 \citep{Bredekamp2009}. The manuscript comprises approximately 240 pages organized into sections distinguished by their illustrations: herbal (botanical), zodiac, cosmological, pharmaceutical, and biological. Despite over a century of scholarly attention, no proposed decipherment has achieved consensus.

Three competing frameworks have dominated Voynich studies. The \emph{cipher hypothesis} holds that the text encodes a natural language through substitution, transposition, or some other cryptographic transformation \citep{dImperio1978, Newbold1928}. The \emph{natural language hypothesis} posits that the script represents a previously unidentified or extinct language, possibly with an unusual writing system \citep{Bax2014}; \citet{Bowern2021} provide a comprehensive linguistic review of this and related approaches. The \emph{hoax hypothesis} contends that the text is meaningless, produced by a mechanical device to simulate linguistic structure without conveying information \citep{Rugg2004, Schinner2007}.

Of these, the hoax hypothesis has been most rigorously formalized. \citet{Rugg2004} demonstrated that a Cardan grille---a card with windows that, when placed over a table of syllable fragments, produces plausible-looking word sequences---could generate text exhibiting many statistical properties of Voynichese. His demonstration was influential because it showed that meaningful-looking text could arise from a trivially simple mechanism. However, Rugg's model remained qualitative; he demonstrated feasibility rather than establishing that the Voynich text was in fact produced this way. \citet{Zandbergen2021} extended the grille approach, showing that a more generic table-and-grille method could reproduce the word-length distribution of Voynichese, but did not test whether the generated text's internal dependency structure matched the real manuscript's. Meanwhile, \citet{Timm2020} proposed an alternative generation mechanism based on self-citation---words generated by copying nearby words with small edits---which reproduces many statistical features of Voynichese. Evidence that the text exhibits positional dependencies inconsistent with simple independent generation has accumulated: \citet{Frodi2020} showed that 90.4\% of symbol lags exhibit significant autocorrelation, and \citet{Lindemann2020} demonstrated that the manuscript's conditional character entropy is lower than any natural language, indicating extreme sequential predictability. However, no prior work has formally tested the Cardan grille's independence assumption using inter-slot transition statistics, nor distinguished between competing generation mechanisms on this basis.

We propose a fourth framework: \emph{lossy generation}. Under this model, the manuscript's text was produced by a combinatorial device---not as a hoax, but as a genuine encoding system. The device mapped real content (botanical descriptions, pharmaceutical recipes, or similar knowledge) through a constrained combinatorial mechanism that introduced systematic information loss. The output preserves structural regularities of the input but degrades semantic precision, producing text that is positionally valid but semantically impoverished relative to its source material.

The analogy to modern large language models (LLMs) is instructive. When an LLM generates text, its output is structurally fluent---grammatical, stylistically consistent, locally coherent---but can be semantically unreliable, a phenomenon termed ``hallucination.'' The Voynich generator, we argue, operated similarly: it produced text that obeys the combinatorial grammar of its mechanism (correct slot structure, valid fragment combinations, appropriate transition dependencies) while losing much of the specific content its operator intended to encode. The key difference is that LLM hallucination arises from statistical pattern completion in a neural network, whereas Voynich hallucination arises from the physical constraints of a finite lookup table.

This paper presents five convergent lines of evidence for the lossy generation model, drawn from a four-phase analysis program:
\begin{enumerate}
  \item \textbf{Mechanical structure}: The vocabulary decomposes into seven positional slots with strong inter-slot dependencies, consistent with row-based table lookup and inconsistent with independent-column mechanisms.
  \item \textbf{Default-filler signatures}: Specific fragments exhibit frequency spikes diagnostic of blank-cell defaults in a lookup table.
  \item \textbf{Section recovery}: Default rates independently recover manuscript section boundaries previously identified by information-theoretic methods.
  \item \textbf{Computational falsification}: A row-based synthetic generator dramatically outperforms an independent-slot generator on transition statistics, directly falsifying the Cardan grille mechanism.
  \item \textbf{Content ablation}: Real text correlates with botanical illustration features while synthetic text does not, proving that content---not mechanism---drives the text-illustration relationship.
\end{enumerate}

% ============================================================
% 2. DATA AND METHODS
% ============================================================
\section{Data and Methods}
\label{sec:methods}

\subsection{Transcription corpus}

Our primary corpus is the Takahashi EVA transcription (version IT2a-n) in IVTFF 2.0 format, comprising 37,886 word tokens, 8,148 word types, and 5,645 hapax legomena (69.3\% of vocabulary). Pages are organized by manuscript section following conventional designations: Herbal ($n = 129$ pages), Stars ($n = 25$), Biological ($n = 19$), Pharmaceutical ($n = 16$), Zodiac ($n = 12$), Cosmological ($n = 10$), Astronomical ($n = 8$), and Text-only ($n = 6$), totaling 225 pages. Of the 129 herbal pages, 120 were available for botanical annotation (9 multi-part folios lacked IIIF image mappings).

For cross-validation, we employed the Zandbergen--Landini (ZL) transcription (version ZL3b-n), which provides an independent reading of the same manuscript pages (35,744 tokens, 9,207 types).

\subsection{Seven-slot positional decomposition}
\label{sec:slotmodel}

Positional decomposition of Voynich words has a substantial history. \citet{Stolfi1997} proposed a three-part prefix-midfix-suffix model, later refined into a crust-mantle-core framework. Most recently, \citet{Zattera2022} developed a 12-slot model achieving 86.6\% token coverage, with each slot cataloging specific glyph positions. Our approach differs from Zattera's in aim: where his 12 slots provide a descriptive positional grammar maximizing coverage, our 7 slots are derived from \emph{functional} analysis---specifically, from testing which positional groupings exhibit statistical dependencies consistent with joint generation from a shared row index. The result is fewer slots, each corresponding to a hypothesized column in the generation table, with measurable inter-slot dependencies, default behaviors, and parametric sub-systems that Zattera's descriptive framework does not test for.

Our initial five-slot model (Initial, Prefix, Gallows, Middle, Suffix), built on \citeauthor{Stolfi1997}'s structural observations, captured 79.1\% of tokens. Subsequent analysis revealed that Slot~4 (Middle) was composite, containing 154 apparent fragment types that decompose into three positional sub-slots following a B?E*T? pattern (Bench, E-series, Terminal). The resulting seven-slot model is:

\begin{equation}
  \text{word} = S_1 \cdot S_2 \cdot S_3 \cdot S_{4a} \cdot S_{4b} \cdot S_{4c} \cdot S_5
  \label{eq:slotmodel}
\end{equation}

\noindent where each slot $S_i$ takes a value from a finite inventory including an empty option (Table~\ref{tab:slotinventory}).

\begin{table}[ht]
\centering
\caption{Seven-slot model: fragment inventories and fill rates.}
\label{tab:slotinventory}
\begin{tabular}{llcrl}
\toprule
Slot & Role & Options & Fill Rate & Top fragments \\
\midrule
$S_1$ & Initial  & 5 + empty & 66\% & \texttt{o}, \texttt{qo}, \texttt{s}, \texttt{d}, \texttt{y} \\
$S_2$ & Prefix   & 6 + empty & 34\% & \texttt{ch}, \texttt{l}, \texttt{r}, \texttt{al}, \texttt{sh}, \texttt{ol} \\
$S_3$ & Gallows  & 11 + empty & 48\% & \texttt{k}, \texttt{t}, \texttt{p}, \texttt{cth}, \texttt{ckh}, \texttt{f} \\
$S_{4a}$ & Bench  & 3 + empty & 21\% & \texttt{ch}, \texttt{h}, \texttt{sh} \\
$S_{4b}$ & E-series & 4 + empty & 37\% & \texttt{e}, \texttt{ee}, \texttt{eee}, \texttt{eeee} \\
$S_{4c}$ & Terminal & 2 + empty & 52\% & \texttt{a}, \texttt{o} \\
$S_5$ & Suffix   & 16 + empty & 88\% & \texttt{y}, \texttt{dy}, \texttt{r}, \texttt{l}, \texttt{iin}, \texttt{in} \\
\bottomrule
\end{tabular}
\end{table}

The seven-slot model yields a theoretical independent combinatorial space of 513,360 possible words. The observed vocabulary of 8,148 types occupies 1.6\% of this space, indicating massive structural constraint.

The Slot~4 sub-decomposition achieves 98.2\% coverage: of all tokens matching the five-slot model, 98.2\% conform to the B?E*T? sub-slot pattern. The remaining 1.8\% (553 tokens) violate the expected column order and are classified as anomalies.

\subsection{Statistical methods}

\paragraph{Slot-pair independence.} We test independence between all pairs of slots using Pearson's chi-squared statistic on contingency tables, with cells filtered to require expected counts $\geq 5$ to satisfy the test's distributional assumptions. Effect sizes are reported as Cram\'{e}r's $V$:
\begin{equation}
  V = \sqrt{\frac{\chi^2}{n \cdot \min(r-1, c-1)}}
  \label{eq:cramersv}
\end{equation}
where $n$ is the total count, $r$ is the number of rows, and $c$ the number of columns contributing cells with expected $\geq 5$.

\paragraph{Default-filler detection.} For each slot's most frequent non-empty fragment, we compute the conditional frequency given that adjacent slots are empty versus filled. A default filler is identified by a large positive rate ratio (filled-when-neighbors-empty versus filled-when-neighbors-filled) and a high Cram\'{e}r's $V$ for the associated contingency table.

\paragraph{Section-level comparisons.} Differences in default rates across manuscript sections are tested using the Kruskal--Wallis $H$-test with epsilon-squared effect size. Pairwise comparisons use Mann--Whitney $U$ with Bonferroni correction.

\paragraph{Botanical correlation.} For each binary botanical feature, we compare the mean ch-default rate (defined as the proportion of words with $S_1 = \emptyset$ and $S_2 = \texttt{ch}$) between pages where the feature is present versus absent. Significance is assessed by permutation test (5,000 permutations, two-tailed). The sign test across distinctive features provides a family-wise assessment.

\paragraph{Transition divergence.} Slot-to-slot transition probabilities are compared between real and synthetic corpora using Jensen--Shannon divergence (JSD):
\begin{equation}
  \text{JSD}(P \| Q) = \frac{1}{2} D_{\text{KL}}(P \| M) + \frac{1}{2} D_{\text{KL}}(Q \| M), \quad M = \frac{1}{2}(P + Q)
  \label{eq:jsd}
\end{equation}
JSD is symmetric, bounded in $[0, 1]$, and avoids the singularities of raw KL divergence. Both distributions use Laplace smoothing ($\alpha = 1$) with identical vocabulary sizes.

\paragraph{Content ablation.} For each botanical feature, we compute the difference in mean ch-default rate between present and absent pages for both real text and for each of 100 synthetic corpora matched in page-level token counts. The ablation $p$-value is the fraction of synthetic $|\Delta|$ values equaling or exceeding the real $|\Delta|$.

\paragraph{E-series co-occurrence.} For each ``E-family'' (set of words identical except for $S_{4b}$), we test whether members co-occur on the same pages more than expected by chance, using a permutation test on pairwise page co-occurrence counts (1,000 permutations). Ordering within pages is assessed by Kendall's $\tau$ across position pairs.

\subsection{Botanical annotation}

We annotated all 120 herbal-section pages with 15 binary botanical features (e.g., \texttt{root\_bulbous}, \texttt{leaf\_lobed}, \texttt{flower\_visible}) based on inspection of high-resolution images from the Yale Beinecke Library IIIF service. Initial annotation of 40 pilot pages was followed by extension to all 120 pages. Feature definitions were designed to capture visually unambiguous morphological characteristics.

\subsection{Synthetic generators}
\label{sec:generators}

We constructed two synthetic text generators to test competing mechanistic hypotheses. A third mechanism---the self-citation model of \citet{Timm2020}, where words are generated by copying nearby words with small edits---is not directly tested here but is addressed in the Discussion (Section~\ref{sec:discussion}). Our two generators isolate the specific question of whether slot values are determined independently or jointly:

\begin{itemize}
  \item \textbf{Generator~A (Independent Slots)}: Each of the seven slots is sampled independently from its marginal frequency distribution in the real corpus. This models the Cardan grille mechanism of \citet{Rugg2004}, where each slot/column is selected without regard to others.
  \item \textbf{Generator~B (Row-Based)}: The row signature ($S_1, S_2, S_3, S_5$) is sampled jointly from the observed distribution of row signatures in the real corpus. The three $S_4$ sub-slots are then sampled independently from their marginal distributions. This models a lookup table where selecting a row determines the values of the outer slots simultaneously, while the middle slots provide within-row variation.
\end{itemize}

Both generators produce corpora of 37,886 tokens matched to the real corpus size. For the ablation test, generators produce page-matched token counts across 100 independent replications.

\subsection{Code review and corrections}

All 14 analysis scripts were subjected to systematic code review. Three critical bugs were identified, all of which were \emph{conservative}---they understated the evidence for the hypothesis:

\begin{enumerate}
  \item The E-series co-occurrence test used incorrect counting logic, yielding 34\% significant families; the corrected analysis yields 74.2\%.
  \item The synthetic generator comparison used asymmetric KL divergence with non-standard epsilon smoothing; the corrected analysis uses JSD with proper Laplace smoothing.
  \item The chi-squared independence tests included cells with expected counts $< 5$; the corrected analysis filters these, and Cram\'{e}r's $V$ values \emph{increased} after correction.
\end{enumerate}

All findings reported in this paper use the corrected analyses. The corrections strengthened every affected result.

\subsection{AI-assisted methodology disclosure}

Botanical annotation was performed with assistance from a vision-language AI model (Claude, Anthropic) applied to IIIF scans of manuscript pages, with human review of annotations. Statistical analyses were implemented in Python scripts developed collaboratively with AI assistance. All code, data, and annotations are available in the project repository. The systematic code review (Section~\ref{sec:methods}) was also AI-assisted, identifying three conservative bugs that, once corrected, strengthened all affected results. The AI tools were used for execution, annotation, and quality control, not for hypothesis generation; the research program and all interpretive decisions were human-directed.

% ============================================================
% 3. RESULTS
% ============================================================
\section{Results}
\label{sec:results}

\subsection{Mechanical structure: Seven positional slots with row-level dependencies}
\label{sec:mechanical}

\subsubsection{Slot decomposition}

The seven-slot model (Equation~\ref{eq:slotmodel}) captures the structure of 79.1\% of corpus tokens under strict parsing, rising to 98.2\% coverage of Slot~4 values under the B?E*T? sub-decomposition. The remaining 20.9\% of tokens require a fallback parse; of these, relaxed parsing (Section~\ref{sec:robustness}) recovers an additional 48.6\%, primarily as compound words.

The observed vocabulary (8,148 types) occupies only 1.6\% of the 513,360-combination theoretical space. A Monte Carlo simulation drawing tokens independently from slot marginal frequencies produces approximately 10,815 distinct types for the same corpus size---33\% more than observed. The vocabulary is \emph{more} constrained than independent frequency-weighted selection predicts, requiring structural dependencies beyond marginal frequencies.

\subsubsection{All slot pairs are dependent}

Table~\ref{tab:chisq} reports chi-squared independence tests for all six tested slot pairs, including both adjacent and non-adjacent pairs.

\begin{table}[ht]
\centering
\caption{Chi-squared independence tests for slot pairs ($n = 29{,}973$ tokens matching the five-slot model). All cells with expected count $< 5$ are filtered. All $p$-values $< 10^{-300}$.}
\label{tab:chisq}
\begin{tabular}{lrrrl}
\toprule
Slot Pair & $\chi^2$ & df & Cram\'{e}r's $V$ & Type \\
\midrule
$S_1 \times S_2$ & 13{,}390 & 30 & 0.299 & Adjacent \\
$S_1 \times S_3$ & 14{,}707 & 40 & 0.313 & Non-adjacent \\
$S_1 \times S_5$ & 3{,}283 & 70 & 0.148 & Non-adjacent \\
$S_2 \times S_5$ & 7{,}556 & 90 & 0.205 & Non-adjacent \\
$S_3 \times S_5$ & 4{,}320 & 120 & 0.134 & Semi-adjacent \\
$S_3 \times S_4$ & 10{,}723 & 352 & 0.212 & Adjacent \\
\bottomrule
\end{tabular}
\end{table}

Every slot pair shows overwhelming statistical dependence ($p < 10^{-300}$), with medium-to-large effect sizes (Cram\'{e}r's $V = 0.13$--$0.31$). Crucially, \emph{non-adjacent} pairs---$S_1 \times S_5$, $S_2 \times S_5$---show significant dependencies with $V > 0.1$. Under a Cardan grille or any independent-column mechanism, non-adjacent slot values should be independent conditional on the grille's window positions. The observed non-adjacent dependencies are inconsistent with independent column selection and consistent with a row-based mechanism where a single row index determines multiple slot values simultaneously.

\subsubsection{The ch-default signature}
\label{sec:chdefault}

Among the candidate default fillers tested, the fragment \texttt{ch} in $S_2$ exhibits the strongest mechanical signature (Table~\ref{tab:defaults}).

\begin{table}[ht]
\centering
\caption{Default fragment candidates: conditional frequency when adjacent slots are empty versus filled.}
\label{tab:defaults}
\begin{tabular}{llrrrl}
\toprule
Slot & Fragment & Rate (adj.\ empty) & Rate (adj.\ filled) & Ratio & $V$ \\
\midrule
$S_2$ & \texttt{ch} & 0.391 & 0.028 & 14.0$\times$ & 0.480 \\
$S_2$ & \texttt{ch} (vs $S_3$) & 0.250 & 0.044 & 5.7$\times$ & 0.287 \\
$S_3$ & \texttt{k} & 0.329 & 0.106 & 3.1$\times$ & 0.243 \\
$S_5$ & \texttt{y} & 0.357 & 0.175 & 2.0$\times$ & 0.173 \\
$S_1$ & \texttt{o} & 0.244 & 0.206 & 1.2$\times$ & 0.043 \\
$S_{4c}$ & \texttt{a} & 0.009 & 0.299 & 0.03$\times$ & 0.217 \\
\bottomrule
\end{tabular}
\end{table}

When $S_1$ is empty, \texttt{ch} appears in $S_2$ at a rate of 39.1\%, compared to only 2.8\% when $S_1$ is filled---a 14-fold increase with $V = 0.48$, the largest effect in the entire analysis. This pattern is difficult to explain linguistically but follows naturally from a lookup table: when the leftmost column yields no content (empty $S_1$), the second column defaults to its most common entry (\texttt{ch}), exactly as a blank cell in a table would be read.

Notably, the fragment \texttt{a} in $S_{4c}$ shows \emph{anti-default} behavior: its rate drops to 0.9\% when $S_5$ is empty, versus 29.9\% when $S_5$ is filled. This identifies \texttt{a} as a structural bridge between Slot~4 and Slot~5, not an independent filler---it obligatorily precedes a suffix.

\subsubsection{Row signatures and Slot~4 variation}

Defining a \emph{row signature} as the tuple ($S_1, S_2, S_3, S_5$), we identify 879 unique signatures in the corpus. Of these, 524 (59.6\%) produce multiple word variants through Slot~4 variation alone. The maximum number of variants per row signature is 34, and $S_{4b}$ (E-series) is the most frequently varying sub-slot (79.1\% of multi-variant rows), followed by $S_{4c}$ (74.6\%) and $S_{4a}$ (54.2\%).

All top-50 row signatures by frequency are universal, appearing on 20 or more pages. Page-specific content is encoded not through unique row inventories but through variable row-selection probabilities across pages.

\subsubsection{The E-series as parametric encoding}

Words that are identical except for their $S_{4b}$ value form \emph{E-families}. If the E-series (empty, \texttt{e}, \texttt{ee}, \texttt{eee}, \texttt{eeee}) encodes a graded parameter, family members should co-occur on the same pages at rates exceeding chance.

Of 97 testable E-families (those with two or more members each having frequency $\geq 10$), 72 (74.2\%) show significantly elevated co-occurrence ($p < 0.05$), against 4.9 expected by chance at the 5\% level. The median observed-to-expected co-occurrence ratio is 3.88, meaning E-family members co-occur on the same page nearly four times more often than chance predicts.

Adjacent E-series values show the strongest co-occurrence: \texttt{e}/\texttt{ee} pairs have a median ratio of 3.07 with 89.6\% significance rate. A weak but highly significant ordering effect is present: lower E-values tend to precede higher ones within pages (Kendall's $\tau = 0.050$, $z = 5.51$, $p < 10^{-6}$).

These patterns are consistent with the E-series encoding a graded quantity---perhaps dosage, intensity, or numerosity---applied to a base concept defined by the row signature.\footnote{Non-peer-reviewed speculation about numeric interpretations of the E-series exists in the Voynich community (e.g., proposed mappings to lunar phases). Our contribution is not the observation that \texttt{e}/\texttt{ee}/\texttt{eee} resembles counting---which is visually apparent---but the rigorous demonstration that E-family members co-occur on the same pages at rates 3.9$\times$ above chance, with a statistically significant ordering effect, establishing that the E-series carries page-level semantic content rather than being a generation artifact.}

\subsection{Section recovery from default rates}
\label{sec:bridging}

If the generator mechanism varies across manuscript sections (e.g., through different table partitions or parameter settings), default rates should differ systematically by section. We computed the ch-default rate (proportion of words with $S_1 = \emptyset$ and $S_2 = \texttt{ch}$) and the overall default rate (mean across five slot-level defaults) for each page.

\begin{table}[ht]
\centering
\caption{Default rate variation across manuscript sections. All tests are Kruskal--Wallis with $\text{df} = 7$.}
\label{tab:sections}
\begin{tabular}{lrrl}
\toprule
Default metric & $H$ & $p$ & $\epsilon^2$ \\
\midrule
$S_1$ (\texttt{o})        & 70.71 & $7.6 \times 10^{-14}$ & 0.317 \\
$S_3$ (\texttt{k})        & 49.79 & $1.6 \times 10^{-9}$  & 0.223 \\
$S_{4c}$ (\texttt{a})     & 38.45 & $3.3 \times 10^{-7}$  & 0.172 \\
$S_2$ (\texttt{ch})       & 30.31 & $1.4 \times 10^{-5}$  & 0.136 \\
Overall (mean of 5) & 47.87 & $4.0 \times 10^{-9}$  & 0.215 \\
$S_5$ (\texttt{y})        &  8.21 & 0.315                  & 0.037 \\
\bottomrule
\end{tabular}
\end{table}

Four of five slot-level defaults and the overall default rate differ significantly across sections (Table~\ref{tab:sections}). The strongest discriminator is the $S_1$ default (\texttt{o}): Zodiac pages show a median \texttt{o}-rate of 0.445, nearly triple the Herbal median of 0.160 ($\epsilon^2 = 0.317$). Different sections default in different slots: Zodiac is high on $S_1$, Pharmaceutical on $S_2$, Stars on $S_3$ and $S_{4c}$.

Herbal pages have the lowest overall default rate (median 0.191), consistent with dense botanical content requiring more slots to be filled. Zodiac pages have the highest (median 0.247), consistent with abstract or formulaic content. This finding independently confirms the section-level vocabulary differences reported by \citet{Montemurro2013} using information-theoretic methods, arrived at here through a wholly different analytical lens (mechanical slot-fill patterns rather than semantic clustering).

\subsection{Botanical content correlation}
\label{sec:botanical}

If the generation device encoded real content, pages depicting morphologically distinctive plants should show lower default rates---more slots filled with content fragments. We tested this prediction on all 120 herbal pages annotated with 15 binary botanical features.

\begin{table}[ht]
\centering
\caption{Botanical features significantly correlated with ch-default rate ($n = 120$ herbal pages, permutation test, 5,000 permutations).}
\label{tab:botanical}
\begin{tabular}{lrrrrl}
\toprule
Feature & $n_{\text{pres}}$ & Mean (pres) & Mean (abs) & \% diff & $p$ \\
\midrule
\texttt{root\_bulbous}  & 31  & 0.122 & 0.169 & $-28\%$ & 0.003 \\
\texttt{root\_visible}  & 104 & 0.150 & 0.204 & $-27\%$ & 0.006 \\
\texttt{leaf\_lobed}    & 50  & 0.139 & 0.170 & $-19\%$ & 0.019 \\
\texttt{leaf\_smooth}   & 65  & 0.171 & 0.141 & $+21\%$ & 0.026 \\
\texttt{leaf\_serrated} & 47  & 0.139 & 0.168 & $-17\%$ & 0.039 \\
\bottomrule
\end{tabular}
\end{table}

Five features reach $p < 0.05$ (Table~\ref{tab:botanical}). Pages depicting plants with bulbous roots show a 28\% lower ch-default rate than those without ($p = 0.003$); lobed leaves, 19\% lower ($p = 0.019$); serrated leaves, 17\% lower ($p = 0.039$). The inverse holds for smooth (featureless) leaves, which show 21\% \emph{higher} defaults ($p = 0.026$). The pattern is consistent: morphologically distinctive features predict lower default rates, meaning more slots filled with content.

The result replicated from a 40-page pilot to the full 120-page set, with \texttt{root\_bulbous} ($p = 0.012 \to 0.003$) and \texttt{leaf\_lobed} ($p = 0.006 \to 0.019$) both surviving replication. Across eight distinctive features (excluding generic features like \texttt{root\_visible}), seven of eight show the predicted direction of lower defaults (sign test $p = 0.07$, two-tailed). The effect sizes are modest (17--28\% differences in default rate), as expected for a lossy system where content is partially degraded. The multiple comparisons concern is addressed in Section~\ref{sec:robustness} and, decisively, by the ablation test in Section~\ref{sec:ablation}, which shows that 0 of 100 synthetic corpora match these correlations.

Row-selection profiles (the distribution of row signatures across pages) showed no correlation with botanical features (Mantel test $r = 0.035$, $p = 0.17$ at $n = 120$). Content encoding operates through default-rate modulation, not through page-specific row inventories.

\subsection{Computational falsification of the Cardan grille}
\label{sec:falsification}

To directly test the Cardan grille hypothesis against the row-based model, we constructed Generators~A and B (Section~\ref{sec:generators}) and compared their outputs to the real corpus across 12 statistical metrics.

\begin{table}[ht]
\centering
\caption{Synthetic generator comparison: single-run metrics. Bold indicates the generator closer to the real corpus.}
\label{tab:synthetic}
\begin{tabular}{lrrrl}
\toprule
Metric & Real & GenA & GenB & Winner \\
\midrule
Word entropy ($H_1$, bits) & 10.47 & 11.88 & \textbf{10.99} & GenB \\
Bigram entropy ($H_2$, bits) & 14.81 & 15.19 & \textbf{15.15} & GenB \\
Zipf exponent & 0.904 & \textbf{0.885} & 1.062 & GenA \\
Word length mean & 5.062 & \textbf{4.677} & 4.626 & GenA \\
Word length std & 1.744 & 1.853 & \textbf{1.639} & GenB \\
Vocabulary size & 8{,}148 & 10{,}393 & \textbf{6{,}558} & GenB \\
Hapax legomena & 5{,}645 & \textbf{6{,}182} & 3{,}130 & GenA \\
Hapax proportion & 0.693 & \textbf{0.595} & 0.477 & GenA \\
Type-token ratio & 0.215 & 0.274 & \textbf{0.173} & GenB \\
Default rate & 0.105 & 0.062 & \textbf{0.138} & GenB \\
\midrule
$S_1 \to S_2$ JSD & 0.000 & 0.094 & \textbf{0.011} & GenB \\
$S_3 \to S_5$ JSD & 0.000 & 0.024 & \textbf{0.002} & GenB \\
\bottomrule
\end{tabular}
\end{table}

Generator~B wins on 8 of 12 metrics overall, but the most diagnostic comparison is the transition divergence. On $S_1 \to S_2$ transitions, GenB achieves JSD $= 0.011$ versus GenA's $0.094$---9$\times$ less divergence. On $S_3 \to S_5$ transitions, the gap widens: GenB's JSD is $0.002$ versus GenA's $0.024$, a factor of 15 (Table~\ref{tab:synthetic}).

GenB's near-zero JSD values ($0.011$ and $0.002$ on a $[0,1]$ scale) indicate near-perfect reproduction of the real corpus's transition structure. GenA's values ($0.094$ and $0.024$) reflect substantial divergence---the independent-slot mechanism fails to capture the dependency structure that the row-based mechanism reproduces automatically.

Bootstrap confidence intervals (100 replications) confirm that neither generator fully captures the real corpus: the real values fall outside the 95\% CI for all seven tested metrics in both generators. The remaining gap in GenB's output---particularly its under-generation of vocabulary (6,558 types versus 8,148 real)---points to row-conditioned $S_4$ variation as the missing mechanism. In GenB, the $S_4$ sub-slots are sampled independently of the row signature; in the real generator, $S_4$ values were likely conditioned on the row, expanding the effective vocabulary.

The metrics where GenA outperforms GenB (Zipf exponent, word length mean, hapax count) are second-order effects: GenA's over-generation of vocabulary through combinatorial explosion accidentally produces hapax rates closer to the real value, but through the wrong mechanism.

\textbf{This comparison computationally falsifies the Cardan grille hypothesis.} The independent-column mechanism proposed by \citet{Rugg2004} cannot reproduce the transition dependencies observed in the real manuscript. A row-based mechanism---where a single table-row selection determines multiple slot values simultaneously---is required.

\subsection{Content ablation: Proof of encoded content}
\label{sec:ablation}

The preceding results establish that the text was produced by a row-based generation device, but they do not determine whether the device encoded real content or produced meaningless output. To distinguish these possibilities, we performed an ablation test: does real text correlate with botanical features while synthetic text from the \emph{same} generator does not?

For each of five botanical features, we computed the difference in mean ch-default rate between present and absent pages for the real corpus and for each of 100 synthetic corpora generated by GenB with page-matched token counts.

\begin{table}[ht]
\centering
\caption{Content ablation test: real versus synthetic botanical correlations ($n = 120$ herbal pages, 100 synthetic replications).}
\label{tab:ablation}
\begin{tabular}{lrrrrl}
\toprule
Feature & Real $\Delta$ & Synth mean & Synth SD & $p_{\text{ablation}}$ & Verdict \\
\midrule
\texttt{root\_bulbous}  & $-0.047$ & $-0.001$ & 0.010 & $< 0.0001$ & CONTENT \\
\texttt{root\_visible}  & $-0.054$ & $+0.000$ & 0.011 & $< 0.0001$ & CONTENT \\
\texttt{leaf\_lobed}    & $-0.032$ & $-0.000$ & 0.007 & $< 0.0001$ & CONTENT \\
\texttt{leaf\_smooth}   & $+0.030$ & $+0.000$ & 0.008 & $< 0.0001$ & CONTENT \\
\texttt{leaf\_serrated} & $-0.029$ & $-0.000$ & 0.008 & $< 0.0001$ & CONTENT \\
\bottomrule
\end{tabular}
\end{table}

All five features yield $p_{\text{ablation}} < 0.0001$ (Table~\ref{tab:ablation}). The real botanical correlations are 3--5 standard deviations from the synthetic null distribution. Zero of 100 synthetic corpora produced differences as extreme as the real corpus for any feature.

This is the decisive result. The generator mechanism alone---even the correct row-based mechanism---cannot produce the observed text-illustration correlations. Real botanical content entered the generation system and survived, degraded but detectable, in the output text. The Voynich Manuscript is not a hoax; it encodes real information through a lossy combinatorial process.

% ============================================================
% 4. ROBUSTNESS CHECKS
% ============================================================
\section{Robustness Checks}
\label{sec:robustness}

\subsection{Cross-transcription validation}

The seven-slot model was developed on the Takahashi (IT) transcription. To test transcription dependence, we examined whether words that fail the slot model (``fallback-hapax'' words) are artifacts of transcription errors by comparing with the Zandbergen--Landini (ZL) transcription.

Of 4,044 IT fallback-hapax types, 62.3\% have the identical reading in ZL at the corresponding locus. Only 15.6\% could be ``rescued'' (converted to a valid five-slot parse) by adopting the ZL reading. Even among words where both transcriptions agree, the hapax-fallback rate remains 68\%. The fallback-hapax phenomenon is overwhelmingly real generator behavior, not transcription noise.

\subsection{Relaxed parser}

A relaxed parser tested four boundary-loosening strategies against the 4,044 fallback-hapax types:

\begin{table}[ht]
\centering
\caption{Relaxed parsing of fallback-hapax words.}
\label{tab:relaxed}
\begin{tabular}{lrr}
\toprule
Strategy & Captured & \% of fallback hapax \\
\midrule
Suffix extension (compound words) & 1{,}690 & 41.8\% \\
d-connector (interstitial \texttt{d})      & 498    & 12.3\% \\
Merged onset ($S_2$+$S_3$ fusion)    & 192    & 4.7\% \\
Expanded $S_1$ (absorb consonant) & 156    & 3.9\% \\
\midrule
Combined (all strategies)          & 1{,}965 & 48.6\% \\
Still uncaptured                   & 2{,}079 & 51.4\% \\
\bottomrule
\end{tabular}
\end{table}

Suffix extension dominates, capturing 41.8\% of fallback hapax as compound words (two slot-sequences concatenated). This suggests the generation device sometimes performed double reads across the table. The 51.4\% that remain uncaptured contain rare EVA glyphs (\texttt{g}, \texttt{x}, \texttt{v}), non-standard initials, or genuinely anomalous glyph combinations---likely a mixture of transcription noise, damaged text, and true generation failures.

\subsection{Anomaly distribution}

The 25.8\% of tokens classified as anomalous (hapax, fallback-only, Slot~4 violations, repetition artifacts) are not uniformly distributed. All anomaly types cluster significantly by page (chi-squared $p < 0.001$). The most anomalous section is Astronomical (47.3\% anomalous), and the least is Biological (16.2\%). Longer pages tend to have lower anomaly rates ($r = -0.23$), consistent with more regular text production during sustained encoding sessions.

\subsection{Code review}

As detailed in Section~\ref{sec:methods}, all three critical bugs discovered during code review were conservative. The corrected E-series co-occurrence rate (74.2\%) is more than double the original (34\%). The corrected JSD values provide properly bounded divergence measures while preserving the 9--15$\times$ GenB advantage. The corrected chi-squared tests show increased Cram\'{e}r's $V$ values. No finding was weakened by correction.

\subsection{Multiple comparisons}

The botanical correlation analysis tests 15 features independently, raising a multiple comparisons concern. We address this in three ways. First, \texttt{root\_bulbous} at $p = 0.003$ survives Bonferroni correction for 15 tests ($\alpha_{\text{adj}} = 0.0033$). Second, the sign test across distinctive features (7/8 in the predicted direction, $p = 0.07$) provides a family-wise assessment that does not depend on individual $p$-values. Third, and most importantly, the ablation test (Section~\ref{sec:ablation}) is immune to multiple comparisons: even if the original $p$-values were inflated, the fact that 0 of 100 synthetic corpora matched the real correlations for \emph{any} feature confirms the result.

% ============================================================
% 5. DISCUSSION
% ============================================================
\section{Discussion}
\label{sec:discussion}

\subsection{The generation device}

Our results constrain the physical form of the generation mechanism. The 879 observed row signatures, the 60\% rate of within-row Slot~4 variation, and the universal distribution of high-frequency signatures across pages are consistent with a lookup table with 879 observed row signatures organized in 7 columns. However, the frequency distribution of row signatures is highly skewed: the top 50 signatures account for a disproportionate share of tokens, and 355 signatures produce only a single word variant. The effective number of ``active'' rows---those contributing meaningfully to the productive vocabulary---is likely in the range of 200--400, with the remainder representing the long tail of specialized or rarely-needed entries. The row determines the outer-slot values ($S_1, S_2, S_3, S_5$); the inner slots ($S_{4a}, S_{4b}, S_{4c}$) provide parametric variation within each row. Common table entries (the top~20 row signatures account for a large share of tokens) serve as structural vocabulary, while the long tail of rare signatures carries more specialized content.

This design is more complex than a Cardan grille but far simpler than a cipher system. The operator would scan the table, select a row corresponding to the concept to be encoded, and read off the word---with the possibility of choosing different $S_4$ values to modify the concept (e.g., different E-series values for different quantities). The physical realization could be a table on paper or parchment, a set of concentric wheels with aligned windows, or a similar device.

\subsection{The E-series and parametric encoding}

The E-series sub-slot ($S_{4b}$: empty, \texttt{e}, \texttt{ee}, \texttt{eee}, \texttt{eeee}) is the most intriguing feature of the generation system. The 74.2\% co-occurrence rate, the 3.88$\times$ median co-occurrence ratio, and the weak but significant ordering effect ($\tau = 0.050$) collectively support the interpretation of the E-series as a graded parameter applied to a base concept.

In a botanical-pharmaceutical context, plausible candidates include dosage (none, small, medium, large), quantity or number (of leaves, petals, roots), intensity (mild, moderate, strong), or temporal duration (brief, extended). The adjacent-value co-occurrence pattern (\texttt{e}/\texttt{ee} co-occurring more strongly than \texttt{e}/\texttt{eee}) is consistent with a scale where nearby values are more likely to co-occur in descriptions of the same subject.

\subsection{Fallback words as generation artifacts}

The 10.2\% of tokens classified as fallback-only (failing the seven-slot parse) and the 14.9\% classified as hapax are, in our framework, the generation system's analog of ``hallucination.'' These words arise when the encoding process produces glyph combinations outside the normal combinatorial grammar---perhaps through misaligned table reads, compound encodings (the 41.8\% suffix-extension rate supports this), or attempts to encode concepts for which the table has no direct entry.

The cross-transcription analysis confirms that these are overwhelmingly real (62.3\% identical in both transcriptions), not artifacts of reading errors.

\subsection{Section-level variation}

Different manuscript sections default in different slots: Zodiac pages show high $S_1$ (\texttt{o}) default rates, Pharmaceutical pages show high $S_2$ (\texttt{ch}) rates, Stars pages show high $S_3$ (\texttt{k}) and $S_{4c}$ (\texttt{a}) rates. This pattern suggests either that different sections used different regions of the generation table, that the table's parameters were adjusted for different content types, or that the nature of the source material varies systematically across sections.

The finding that Herbal pages have the lowest overall default rate (median 0.191) while Zodiac pages have the highest (0.247) is consistent with botanical descriptions being the table's primary design target---the section for which the encoding device was most expressive.

\subsection{Relationship to Timm \& Schinner's self-citation model}

\citet{Timm2020} proposed that Voynich words are generated by copying nearby words with small character-level edits, a mechanism that reproduces many statistical features of Voynichese including its characteristic word-length distribution and high local repetition rate. Our row-based model is compatible with, but distinct from, self-citation. In a table-lookup framework, the same rows tend to be reused on nearby pages because the author was encoding related content sequentially---producing the same local similarity that the self-citation model captures. However, self-citation alone does not predict the slot-dependency structure ($V = 0.12$--$0.31$ across non-adjacent pairs), the ch-default signature ($V = 0.48$), or the E-series parametric co-occurrence (74.2\% of families). These features require a structured generation mechanism, not merely local copying with noise. The page-proximity effect we observe ($r = 0.20$, $p = 0.015$ for row-profile similarity between adjacent folios) is consistent with both models: the author reused similar table regions for nearby pages, which would produce the same statistical signature as self-citation. Our framework explains \emph{why} self-citation patterns emerge---because the same table rows get reused on nearby pages---while providing additional explanatory power for the dependency and default structure.

\subsection{Implications for other manuscript features}

The page-proximity effect---adjacent folios tend to share similar statistical profiles---may reflect the encoding process (the operator working through a sequence of related subjects) or may provide evidence for folio reordering. The high anomaly rates in Astronomical (47.3\%) and Zodiac (44.5\%) sections may indicate that these sections stretch the generation device beyond its design parameters, attempting to encode concepts for which the table has fewer dedicated rows.

\subsection{Limitations}

Several limitations constrain our conclusions. First, botanical annotation was AI-assisted, and while annotations were reviewed, systematic biases in the vision model could affect the correlation analysis. The ablation test mitigates this concern: any annotation bias would be shared between real and synthetic analyses, so it cannot produce the observed ablation result.

Second, the Slot~4 sub-decomposition, while achieving 98.2\% coverage, leaves the 1.8\% column-order violations unexplained. These may represent alternative access patterns, scribal errors, or places where our decomposition model is inexact.

Third, the botanical effect sizes are modest (17--28\% differences in default rate), and the sign test across distinctive features is marginal ($p = 0.07$). The lossy generation framework predicts modest effects: real content passes through the system but is substantially degraded. The ablation test's decisive result ($p < 0.0001$ for all five features) compensates for the modest individual effect sizes by demonstrating that the effects, however small, are genuine content signals rather than mechanical artifacts.

Fourth, our model of GenB treats $S_4$ sub-slots as independent of the row signature, but the vocabulary gap (6,558 vs.\ 8,148 types) indicates that the real device conditioned $S_4$ on the row. A fuller model incorporating row-$S_4$ coupling would likely close this gap but requires assumptions about the coupling structure that we cannot currently constrain.

% ============================================================
% 6. CONCLUSION
% ============================================================
\section{Conclusion}
\label{sec:conclusion}

We have presented five convergent lines of evidence that the Voynich Manuscript's text was produced by a lossy combinatorial generation device.

\begin{enumerate}
  \item The text decomposes into seven positional slots with strong inter-slot dependencies (all pairs $p < 10^{-300}$, $V = 0.12$--$0.31$), consistent with row-based table lookup.
  \item The fragment \texttt{ch} in Slot~2 shows a 14$\times$ default-filler signature ($V = 0.48$), diagnostic of blank-cell defaults in a lookup table.
  \item Default rates recover manuscript section boundaries ($p = 4 \times 10^{-9}$), confirming prior information-theoretic findings via a wholly different method.
  \item A row-based generator reproduces real transition statistics with 9--15$\times$ less divergence than an independent-slot generator, computationally falsifying the Cardan grille model of \citet{Rugg2004}.
  \item Real text correlates with botanical features ($p = 0.003$ for \texttt{root\_bulbous}) while synthetic text does not ($p_{\text{ablation}} < 0.0001$), proving that genuine content passed through the generation system.
\end{enumerate}

The Voynich Manuscript is neither a simple cipher, a natural language transcription, nor a meaningless hoax. It is the output of a mechanical encoding system that mapped real botanical and pharmaceutical knowledge through a constrained combinatorial process, producing text that preserves structural regularities while losing semantic precision. The ch-default rate emerges as a new analytical tool for Voynich studies, sensitive to both section-level structure and page-level content density.

The lossy generation framework has implications beyond the Voynich Manuscript. For any suspected generated text---historical or modern---the analytical program is: decompose into positional units, test for inter-unit dependencies, identify default signatures, construct competing synthetic generators, and compare real versus synthetic correlations with ground truth. This methodology can distinguish mechanical generation from natural composition, and content-bearing generation from empty fabrication. If our analysis is correct, the generation artifacts in the Voynich Manuscript---the fallback words, the default-filled tokens, the combinatorial hallucinations---represent the first documented text-generation artifacts, predating computational text generation by at least five centuries.

% ============================================================
% ACKNOWLEDGMENTS
% ============================================================
\section*{Acknowledgments}

We thank the Yale Beinecke Rare Book and Manuscript Library for providing IIIF access to high-resolution scans of MS 408. The Takahashi transcription (IT2a-n) and Zandbergen--Landini transcription (ZL3b-n) are indispensable community resources without which quantitative Voynich research would be impossible. AI tools (Claude, Anthropic) were used for coding assistance, botanical annotation, code review, and manuscript review as disclosed in Section~\ref{sec:methods}.

% ============================================================
% REFERENCES
% ============================================================
\begin{thebibliography}{99}

\bibitem[Bax, 2014]{Bax2014}
Bax, S. (2014).
\newblock A proposed partial decoding of the Voynich Manuscript.
\newblock {\em CAS Working Papers on Voynich Studies}.

\bibitem[Bowern and Lindemann, 2021]{Bowern2021}
Bowern, C. and Lindemann, L. (2021).
\newblock The linguistics of the Voynich Manuscript.
\newblock {\em Annual Review of Linguistics}, 7:285--308.

\bibitem[Bredekamp et~al., 2009]{Bredekamp2009}
Bredekamp, H., et~al. (2009).
\newblock Radiocarbon dating of the Voynich Manuscript.
\newblock University of Arizona Accelerator Mass Spectrometry Laboratory. Reported in multiple sources; vellum dated to 1404--1438 CE.

\bibitem[Frod\'{i} et~al., 2020]{Frodi2020}
Frod\'{i}, T., et~al. (2020).
\newblock The Voynich Manuscript: Symbol roles revisited.
\newblock {\em PLOS ONE}, 15(8):e0237066.

\bibitem[Gaskell and Bowern, 2022]{Gaskell2022}
Gaskell, D. and Bowern, C. (2022).
\newblock Gibberish after all? Voynich Manuscript vs.\ human-generated gibberish.
\newblock In {\em CEUR Workshop Proceedings}, volume 3313.

\bibitem[Currier, 1976]{Currier1976}
Currier, P. (1976).
\newblock Papers on the Voynich Manuscript.
\newblock Unpublished manuscript, New Elizabethan Cipher Symposium.

\bibitem[d'Imperio, 1978]{dImperio1978}
d'Imperio, M.~E. (1978).
\newblock {\em The Voynich Manuscript: An Elegant Enigma}.
\newblock National Security Agency, Fort Meade, MD.

\bibitem[Lindemann and Bowern, 2020]{Lindemann2020}
Lindemann, L. and Bowern, C. (2020).
\newblock Character entropy in modern and historical texts: Comparison metrics for an undeciphered manuscript.
\newblock arXiv preprint arXiv:2010.14697.

\bibitem[Landini, 2001]{Landini2001}
Landini, G. (2001).
\newblock Evidence of linguistic structure in the Voynich Manuscript using spectral analysis.
\newblock {\em Cryptologia}, 25(4):275--295.

\bibitem[Montemurro and Zanette, 2013]{Montemurro2013}
Montemurro, M.~A. and Zanette, D.~H. (2013).
\newblock Keywords and co-occurrence patterns in the Voynich Manuscript: An information-theoretic analysis.
\newblock {\em PLoS ONE}, 8(6):e66344.

\bibitem[Newbold, 1928]{Newbold1928}
Newbold, W.~R. (1928).
\newblock {\em The Cipher of Roger Bacon}.
\newblock University of Pennsylvania Press, Philadelphia.

\bibitem[Reddy and Knight, 2011]{Reddy2011}
Reddy, S. and Knight, K. (2011).
\newblock What we know about the Voynich Manuscript.
\newblock In {\em Proceedings of the 5th ACL-HLT Workshop on Language Technology for Cultural Heritage, Social Sciences, and Humanities}, pages 78--86.

\bibitem[Rugg, 2004]{Rugg2004}
Rugg, G. (2004).
\newblock An elegant hoax? A possible solution to the Voynich Manuscript.
\newblock {\em Cryptologia}, 28(1):31--46.

\bibitem[Rugg and Taylor, 2017]{RuggTaylor2017}
Rugg, G. and Taylor, G. (2017).
\newblock Hoaxing statistical features of the Voynich Manuscript.
\newblock {\em Cryptologia}, 41(3):247--268.

\bibitem[Schinner, 2007]{Schinner2007}
Schinner, A. (2007).
\newblock The Voynich Manuscript: Evidence of the hoax hypothesis.
\newblock {\em Cryptologia}, 31(2):95--107.

\bibitem[Sterneck et~al., 2021]{Sterneck2021}
Sterneck, R., Polish, A., and Bowern, C. (2021).
\newblock Topic modeling in the Voynich Manuscript.
\newblock arXiv preprint arXiv:2107.02858.

\bibitem[Stolfi, 1997]{Stolfi1997}
Stolfi, J. (1997).
\newblock A Voynich Manuscript word-structure model.
\newblock Technical report, UNICAMP, Campinas, Brazil.

\bibitem[Timm and Schinner, 2020]{Timm2020}
Timm, T. and Schinner, A. (2020).
\newblock A possible generating algorithm of the Voynich Manuscript.
\newblock {\em Cryptologia}, 44(1):1--19.

\bibitem[Zattera, 2022]{Zattera2022}
Zattera, M. (2022).
\newblock A new transliteration alphabet brings new evidence of word structure and multiple ``languages'' in the Voynich Manuscript.
\newblock In {\em CEUR Workshop Proceedings}, volume 3313.

\bibitem[Zandbergen, 2023]{Zandbergen2023}
Zandbergen, R. (2023).
\newblock The Voynich Manuscript.
\newblock \url{http://www.voynich.nu/}.

\bibitem[Zandbergen, 2021]{Zandbergen2021}
Zandbergen, R. (2021).
\newblock A computational approach to the Voynich Manuscript using table and grille methods.
\newblock arXiv preprint arXiv:2104.10957.

\end{thebibliography}

\end{document}
